\section{Verification and Testing Framework}
\label{sec:verification}

Every result in this report depends on code written for this project: seven
basis functions, six integrators, the training loop and every diagnostic.
Before any of that code is trusted to produce a benchmark number, it is
checked against cases whose correct answer is known independently. This
section describes those checks and the pipeline that turns one training run
into the numbers in the tables.

\subsection{Solver verification against a known solution}

The integrators are checked on $\dot y=-y$, $y(0)=1$, whose exact solution
$y(t)=e^{-t}$ is known. Integrating to $t=1$ with $n$ equal steps gives an
error $e(h)=|y_n-e^{-1}|$ at step $h=1/n$, and by
Equation~\ref{eq:global-error} the observed order between two step sizes is
\begin{equation}
\hat p=\frac{\log e(h_1)-\log e(h_2)}{\log h_1-\log h_2}.
\label{eq:observed-order}
\end{equation}
Table~\ref{tab:observed-order} shows that every solver reaches its
theoretical order. Halving the step divides the error by $2^{p}$, as
predicted. It also shows a difference that order alone does not capture: at
the same fifth order, Tsit5's error is about $5\times$ smaller than
DOPRI5's, because Tsitouras chose its coefficients to minimize the leading
error constant $K$ in Equation~\ref{eq:global-error}
\cite{tsitouras2011tsit5}.

\begin{table}[!ht]
\centering
\footnotesize
\renewcommand{\arraystretch}{1.15}
\caption{Observed convergence order of this project's integrators on
$\dot y=-y$ over $[0,1]$ (double precision, one substep). The automated
test requires $|\hat p-p|$ below $0.20$ (Euler) to $0.45$ (fifth order).}
\label{tab:observed-order}
\begin{tabular}{@{}L{2.4cm} C{1.2cm} C{2.0cm} C{2.0cm} C{2.0cm} C{1.6cm} C{1.6cm}@{}}
\toprule
\hdrrow \thh{Solver} & \thh{Order $p$} & \thh{$e(h{=}0.1)$} & \thh{$e(h{=}0.05)$} & \thh{$e(h{=}0.025)$} & \thh{$\hat p$ (0.1$\to$0.05)} & \thh{$\hat p$ (0.05$\to$0.025)}\\
\midrule
Forward Euler & 1 & $1.92\times10^{-2}$ & $9.39\times10^{-3}$ & $4.65\times10^{-3}$ & 1.03 & 1.02\\
Midpoint & 2 & $6.62\times10^{-4}$ & $1.59\times10^{-4}$ & $3.91\times10^{-5}$ & 2.06 & 2.03\\
Heun & 2 & $6.62\times10^{-4}$ & $1.59\times10^{-4}$ & $3.91\times10^{-5}$ & 2.06 & 2.03\\
Classical RK4 & 4 & $3.33\times10^{-7}$ & $2.00\times10^{-8}$ & $1.22\times10^{-9}$ & 4.06 & 4.03\\
DOPRI5 & 5 & $1.21\times10^{-9}$ & $3.48\times10^{-11}$ & $1.04\times10^{-12}$ & 5.12 & 5.06\\
Tsit5 & 5 & $2.56\times10^{-10}$ & $6.43\times10^{-12}$ & $1.77\times10^{-13}$ & 5.31 & 5.18\\
\bottomrule
\end{tabular}
\end{table}

Midpoint and Heun agree to every digit here because, for a linear field,
both second-order methods reduce to the same update
$y_{n+1}=(1-h+h^2/2)\,y_n$. They differ only on nonlinear fields such as the
learned ones. Further solver tests check that the Tsit5 and DOPRI5 tableaus
are internally consistent (weights summing to one, the embedded fourth-order
pair consistent), that energy on a harmonic oscillator drifts orders of
magnitude less under RK4 and Tsit5 than under Euler, that time-dependent and
backward-in-time integration are correct, and that gradients flow through
every solver.

\subsection{The automated test suite}

The checks run as an automated \texttt{pytest} suite of 250 tests. On the
final code, 249 pass and 1 is skipped: that test measures peak GPU memory
and runs only when a CUDA device is present. Table~\ref{tab:test-suite}
groups the tests by what they protect.

\begin{table}[!ht]
\centering
\footnotesize
\renewcommand{\arraystretch}{1.15}
\caption{Automated tests, grouped by the component they verify.}
\label{tab:test-suite}
\begin{tabularx}{\textwidth}{@{}L{3.6cm} C{1.1cm} X@{}}
\toprule
\hdrrow \thh{Component} & \thh{Tests} & \thh{What is checked}\\
\midrule
Integrators & 32 & observed order (Table~\ref{tab:observed-order}); tableau consistency; energy drift; substep monotonicity; non-autonomous and backward integration; batching; gradient flow\\
Basis functions and KAN layer & 57 & shape and validity of all seven bases; exact known values (Chebyshev values and recurrence, Newton values, Lagrange cardinal property); partition of unity (B-spline, Lagrange); gradient flow through every basis; layer forward and backward passes\\
MLP baseline & 22 & parameter count and KAN/MLP parameter matching; supported activations; initialization statistics; gradient flow through the ODE; seed reproducibility\\
Datasets & 10 & shapes, splits and reproducibility; pendulum energy conserved at $\mu=0$ and dissipated at $\mu>0$; the SIR invariant $S+I+R=1$; epidemic-curve normalization\\
Metrics & 11 & MSE, MAE, $R^2$, relative $L_2$, gradient norm and function-evaluation counts on inputs with known answers; Lipschitz estimate on a linear map and on $\sin$\\
Training pipeline & 23 & end-to-end training; the selected checkpoint is the best one, not the last; a reloaded checkpoint rebuilds the same field, including every fix; rescaling time equals rescaling the field, and $s=1$ changes nothing bit for bit; the projected field sums to zero far outside the training window; the gated field vanishes on $u_d=0$ (Eqs.~\ref{eq:time-scale}--\ref{eq:vanishing-gate})\\
Plotting & 7 & every figure function runs on real outputs\\
Study-specific checks & 88 & adjoint vs.\ direct gradient agreement (8); gradient-noise statistic (18); hybrid-basis gate (10); sparse-regression and pruning procedures (33); damping-study verdict classifier and non-finite guard (19)\\
\midrule
\textbf{Total} & \textbf{250} & 249 passed, 1 skipped (requires CUDA)\\
\bottomrule
\end{tabularx}
\end{table}

\subsection{From a configuration to a table entry}

Figure~\ref{fig:pipeline} shows the path every benchmark number takes. One
function trains any configuration (system, model, basis, solver, epochs,
seed and fixes) and writes a self-describing results record: the full
configuration, the software version and code revision, and every metric at
both the selected and the final checkpoint. A benchmark driver runs lists of
configurations through that function, and the tables in this report are
collated from the saved records rather than typed by hand, so any number can
be traced back to the run that produced it.

\paragraph{Post-hoc analyses.} A few numbers were computed afterwards from
saved checkpoints and training records, without training anything: the
pendulum models rescored on a common test window, the hybrid basis's
measured contributions, the RBF/B-spline comparison at a matched budget, the
far-horizon scores of the 50{,}000-epoch models, the Newton root-finding and
Jacobian spectra of the damping study, the damping study's trajectories
redrawn for its figures, and the exact-derivative sparse regression. Each is
produced by a short script in the project repository
(\texttt{implementation/experiments/posthoc\_checks}), and every redrawn
report figure by a script in \texttt{report/figure\_scripts}. These scripts
were added with this revision of the report, after the results commit named
on the title page, and are not covered by the automated tests of
Table~\ref{tab:test-suite}. Their outputs were checked against the saved
records they build on instead: each rebuilt checkpoint reproduces its own
recorded metrics (for example, the 10{,}000-epoch far-horizon scores to
three significant figures, the pendulum and SIR $R^2$ exactly, and every
damping-study trajectory its recorded best training MSE), and the
finite-difference sparse regression reproduces the saved noise-free result.

\begin{figure}[htbp]
\centering
\resizebox{\textwidth}{!}{%
\begin{tikzpicture}[font=\small,
  st/.style={draw=ink!60, fill=codebg, rounded corners=2pt, minimum height=9mm, text width=2.35cm, align=center, font=\scriptsize},
  key/.style={draw=accent, fill=accentsoft, rounded corners=2pt, minimum height=9mm, text width=2.2cm, align=center, font=\scriptsize},
  arr/.style={-{Stealth[length=2mm]}, draw=ink!65, thick}]
  \node[st]  (cfg) at (0,0)     {Configuration\\system, model, basis,\\solver, epochs, seed};
  \node[st]  (dat) at (2.9,0)   {Data generator\\reference solve,\\train/test split};
  \node[key] (trn) at (5.8,0)   {Training loop\\Algorithm~\ref{alg:train}};
  \node[st]  (ckp) at (8.7,0)   {Checkpoint with\\lowest\\training loss};
  \node[st]  (scr) at (11.6,0)  {Scoring on the\\full horizon};
  \node[st]  (rec) at (14.5,0)  {Results record\\config, metrics,\\history, figures};
  \node[st]  (tab) at (14.5,-2.0) {Collated tables\\and report figures};
  \node[key] (tst) at (5.8,-2.0) {Automated tests\\(Table~\ref{tab:test-suite})};
  \draw[arr] (cfg) -- (dat); \draw[arr] (dat) -- (trn); \draw[arr] (trn) -- (ckp);
  \draw[arr] (ckp) -- (scr); \draw[arr] (scr) -- (rec); \draw[arr] (rec) -- (tab);
  \draw[arr, dashed, accent] (tst) -- node[right, font=\tiny\sffamily\color{accent}]{verify every component first} (trn);
  \node[font=\tiny\sffamily\color{ink!60}] at (8.7,-0.75) {test window never used here};
\end{tikzpicture}}
\caption{The benchmark pipeline. Model selection uses only the training
window; the test window is used once, when the selected checkpoint is
scored. The same code path produces every table in this report.}
\label{fig:pipeline}
\end{figure}

\subsection{Metrics}

A model's predicted trajectory $\hat{\mathbf u}$ is compared with the
reference $\mathbf u$ on a set $\mathcal T$ of time points (the test window
unless stated otherwise), with $d$ state components:
\begin{align}
\mathrm{MSE}&=\frac{1}{|\mathcal T|\,d}\sum_{t\in\mathcal T}\lVert\hat{\mathbf u}(t)-\mathbf u(t)\rVert_2^2,
\qquad \mathrm{RMSE}=\sqrt{\mathrm{MSE}},\label{eq:metric-mse}\\
R^2&=1-\frac{\sum_{t\in\mathcal T}\lVert\hat{\mathbf u}(t)-\mathbf u(t)\rVert_2^2}{\sum_{t\in\mathcal T}\lVert\mathbf u(t)-\bar{\mathbf u}\rVert_2^2},
\qquad
\text{rel.\ }L_2=\frac{\lVert\hat{\mathbf u}-\mathbf u\rVert_2}{\lVert\mathbf u\rVert_2},\label{eq:metric-r2}
\end{align}
where $\bar{\mathbf u}$ is the per-component mean over $\mathcal T$. $R^2=1$
is a perfect match, $R^2=0$ is no better than predicting that mean, and a
negative $R^2$ is worse than it. The smoothness of a learned field is
summarized by an empirical Lipschitz estimate,
\begin{equation}
\hat L=\max_{k}\frac{\lVert\mathbf f_\theta(\mathbf u_k+\boldsymbol\epsilon_k)-\mathbf f_\theta(\mathbf u_k)\rVert_2}{\lVert\boldsymbol\epsilon_k\rVert_2},
\qquad \lVert\boldsymbol\epsilon_k\rVert_2=10^{-3},
\label{eq:lipschitz}
\end{equation}
taken over the training states $\mathbf u_k$, each with a random perturbation
direction. This is the largest local slope of the field observed along the
training trajectory. It is a \emph{lower} bound on the field's true Lipschitz
constant, which would require a maximum over all pairs of points, and it is
used only to compare how sharp different trained fields are. Training cost
is reported as function evaluations per epoch (Equation~\ref{eq:nfe}) and as
wall-clock time. All timings are CPU timings (a local development machine,
or cloud CPUs for the extended-budget runs of Section~\ref{sec:phase4}),
except the adjoint-versus-autograd comparison of Section~\ref{sec:track-a},
which measures GPU memory and runs on an RTX~3060. Timings are compared
only within one machine, except in the cost note of Section~\ref{sec:phase4},
where the difference between machines is itself the point.

\FloatBarrier
